% =====================================================================
%  Chapter 4, Part 1: Conception
%  Sources: docs/adr/0001-0012 of the code repository.
% =====================================================================

% =====================================================================
\section*{Introduction}
\addcontentsline{toc}{section}{Introduction}
\phantomsection\label{sec-d-intro}
% =====================================================================

This chapter presents the design and the development of the framework
whose contributions were described in \cref{chap-contribution}. The
framework was built in thirteen short iterative sprints. Each sprint
delivered a working part of the system (the grounded tutor, the adaptive
loop, the context model, the teacher pages, the in-app reader, the
Material Improvement agent, the learning-gain evaluation, explainable
personalization, the LTI launch and the pilot kit, the early warning), and
each major decision was recorded in an \emph{architecture decision record}
(ADR): a short document that states the context, the decision, the
alternatives and the consequences~\cite{nygard2011adr}. Twelve ADRs were
written; they are the source of the design choices presented here.

The first part of the chapter covers the design. It analyses the
requirements (\cref{sec-d-req}), then presents the architecture
(\cref{sec-d-arch}), the dynamic view of the main scenarios
(\cref{sec-d-dynamic}), the static view of the domain
(\cref{sec-d-static}), the design of the database (\cref{sec-d-db}) and
the technology stack (\cref{sec-d-stack}). The second part covers the
development: the process, the code base, the interfaces, the tests and
the deployment.

% =====================================================================
\section{Requirements analysis}
\label{sec-d-req}
% =====================================================================

\subsection{Actors}

\Cref{tab-actors} lists the actors of the system. The two primary actors
are the people who use the platform; the three secondary actors are
systems that the platform depends on.

\begin{table}[htbp]
\centering
\caption{Actors of the system}
\label{tab-actors}
\small
\begin{tabularx}{\textwidth}{@{}L{2.6cm}L{2.0cm}Y@{}}
\toprule
\textbf{Actor} & \textbf{Type} & \textbf{Role} \\
\midrule
Student & primary & reads the course, asks the tutor, takes quizzes, sets a goal, follows the path and the suggestions, takes part in a study \\
Teacher & primary & indexes the course, reviews generated questions, sets objectives, steers recommendations, monitors the class, acts on early warnings, improves the material, runs a study \\
Moodle & secondary & provides courses, activities, files, enrolments, completion and grades; launches the app (LTI 1.3) and receives the progress grade (AGS) \\
LLM provider & secondary & computes embeddings and writes answers, quiz items and material drafts (hosted or local) \\
Scheduler & secondary & starts the periodic context sync, after which every student's risk is scored (BullMQ repeatable job) \\
\bottomrule
\end{tabularx}
\end{table}

\subsection{Functional requirements}

\Cref{tab-fr} lists the functional requirements, numbered by actor:
FR-S for the student, FR-T for the teacher and FR-Y for the system.

\begin{table}[htbp]
\centering
\caption{Functional requirements}
\label{tab-fr}
\small
\begin{tabularx}{\textwidth}{@{}L{1.4cm}Y@{}}
\toprule
\textbf{Id} & \textbf{Requirement} \\
\midrule
\multicolumn{2}{@{}l}{\textit{Student}} \\
FR-S1 & Sign in from Moodle (LTI 1.3 launch), without a second login. \\
FR-S2 & See my courses and my progress. \\
FR-S3 & Read the course activities inside the app (pages, PDFs) without a second login. \\
FR-S4 & Ask the tutor and see its steps and its sources. \\
FR-S5 & Get hints, not solutions, for exercises. \\
FR-S6 & Take adaptive quizzes with immediate feedback. \\
FR-S7 & See my mastery per section, my deadlines and my misconceptions. \\
FR-S8 & Receive proactive suggestions with their reason, and dismiss them. \\
FR-S9 & Set my goal, my level and my preferred explanation style for a course. \\
FR-S10 & See my learning path and why each section is recommended. \\
FR-S11 & Rate the tutor's answers. \\
FR-S12 & Take part in a study: consent or decline, take the pre-test, the post-test and the SUS questionnaire, withdraw. \\
\midrule
\multicolumn{2}{@{}l}{\textit{Teacher}} \\
FR-T1 & Index the course material. \\
FR-T2 & Generate, review, edit, approve or reject quiz questions. \\
FR-T3 & Choose the review mode of the course. \\
FR-T4 & Set the target mastery and the key sections. \\
FR-T5 & Monitor the class: mastery heatmap, risk flags, shared misconceptions. \\
FR-T6 & Run the material analysis; accept, edit, reject or mark as applied each suggestion. \\
FR-T7 & Pin or skip a section for one student or for the whole class. \\
FR-T8 & See the early warnings with their reasons; send a nudge, mark a warning handled or dismiss it; see whether the student came back. \\
FR-T9 & Run a study: assign the arms, build the pre-test and post-test, export the pseudonymized results and the pseudonym key. \\
\midrule
\multicolumn{2}{@{}l}{\textit{System}} \\
FR-Y1 & Sync the context periodically. \\
FR-Y2 & Record in-app reading. \\
FR-Y3 & Trace every agent decision. \\
FR-Y4 & Score the risk of every student after each sync and keep a daily snapshot. \\
FR-Y5 & Send each student's course progress to the Moodle gradebook. \\
\bottomrule
\end{tabularx}
\end{table}

\subsection{Non-functional requirements}

The non-functional requirements follow from the contributions and from
the context of an Algerian university.

\begin{description}[leftmargin=1.5cm, labelwidth=1.3cm, labelsep=0.2cm, style=sameline, font=\normalfont\bfseries]
  \item[NFR-1] \textbf{Moodle core untouched.} The integration uses only
        web services and LTI~1.3; no plugin is required, and the only
        write is the progress grade, through LTI's grade service.
  \item[NFR-2] \textbf{Privacy.} A single Moodle token, held by one
        service; signed, short-lived file links; an option to run the
        language model locally; consent, withdrawal and pseudonymized
        exports in a study.
  \item[NFR-3] \textbf{Traceability.} Every routing decision and agent
        step is logged with a trace identifier.
  \item[NFR-4] \textbf{Grounding.} Answers come only from the course
        material, with citations, or the tutor says that the answer is not
        in the course.
  \item[NFR-5] \textbf{Responsiveness.} Answers are streamed; long tasks
        (ingestion, quiz writing, material drafts) run in the background.
  \item[NFR-6] \textbf{Testability.} LLMs are replaced by fake models in
        the automated tests, so the tests never call an LLM.
  \item[NFR-7] \textbf{Bilingual interface} (French and English).
  \item[NFR-8] \textbf{Deployability.} The whole stack starts with Docker
        Compose.
  \item[NFR-9] \textbf{Explainability.} Every recommendation and every
        warning carries the structured reasons of the rule or model that
        produced it; no explanation is written by an LLM.
\end{description}

\subsection{Use cases}

\Cref{fig-use-cases} shows the use case diagram. The student and the
teacher each have five use cases. Asking the tutor always includes the
retrieval of course passages, and taking an adaptive quiz always includes
the update of the mastery. The Scheduler starts the context sync.

\begin{figure}[htbp]
  \centering
  \reportfig{use-cases}
  \caption{Use case diagram.}
  \label{fig-use-cases}
\end{figure}

\Crefrange{tab-uc-tutor}{tab-uc-warning} describe the five main use
cases.

\newcommand{\ucrow}[2]{\textbf{#1} & #2 \\}
\begin{table}[htbp]
\centering
\caption{Use case \emph{Ask the tutor}}
\label{tab-uc-tutor}
\small
\begin{tabularx}{\textwidth}{@{}L{2.6cm}Y@{}}
\toprule
\ucrow{Actors}{Student (primary); LLM provider (secondary)}
\ucrow{Pre-conditions}{The student is signed in and enrolled; the course is indexed.}
\ucrow{Main scenario}{1. The student writes a question in the course chat.
  2. The Orchestrator classifies it (intent, language, topic).
  3. The Tutor retrieves the closest passages of the course.
  4. The best similarity reaches the threshold; the Tutor adapts the level
  to the mastery and context of the student.
  5. The LLM writes the answer, streamed with the agent steps.
  6. The citations are checked and shown with the answer.}
\ucrow{Alternatives}{2a. Smalltalk: a short reply ends the turn.
  2b. Quiz intent: the tutor proposes a quiz on the section.
  4a. The best similarity is below the threshold: the tutor says that the
  answer is not in the course and lists the closest sections.
  4b. The question is an exercise: the tutor gives hints only.}
\ucrow{Post-conditions}{The messages, steps, citations and a context snapshot are stored in one transaction; the turn is traced.}
\bottomrule
\end{tabularx}
\end{table}

\begin{table}[htbp]
\centering
\caption{Use case \emph{Take an adaptive quiz}}
\label{tab-uc-quiz}
\small
\begin{tabularx}{\textwidth}{@{}L{2.6cm}Y@{}}
\toprule
\ucrow{Actors}{Student (primary); LLM provider (secondary)}
\ucrow{Pre-conditions}{The student is enrolled; the section has ingested material.}
\ucrow{Main scenario}{1. The student starts a quiz on a section.
  2. The Assessment agent fills the quiz from approved bank questions that
  the student has never answered.
  3. The quiz is ready; items are shown without their keys.
  4. For each answer, the system grades it in code, updates the mastery
  (BKT) and records the misconception of a wrong option.
  5. The student sees the feedback, the explanation and the mastery before
  and after.
  6. The Recommendation agent proposes the next section with its reason.}
\ucrow{Alternatives}{2a. Shortfall in \texttt{auto\_publish} mode: the LLM writes the missing items, they are validated in code and used.
  2b. Shortfall in \texttt{review} mode: the missing items are written as pending; the quiz is served if enough questions are approved, otherwise it ends with \texttt{awaiting\_review}.
  2c. The LLM fails or times out: the quiz ends as \texttt{failed} with a code that the interface explains.}
\ucrow{Post-conditions}{Responses and the new learner state are committed together; the mastery changes are traced.}
\bottomrule
\end{tabularx}
\end{table}

\begin{table}[htbp]
\centering
\caption{Use case \emph{Review the question bank}}
\label{tab-uc-bank}
\small
\begin{tabularx}{\textwidth}{@{}L{2.6cm}Y@{}}
\toprule
\ucrow{Actors}{Teacher (primary); LLM provider (secondary)}
\ucrow{Pre-conditions}{The teacher is signed in and teaches the course.}
\ucrow{Main scenario}{1. The teacher opens the question bank of a section.
  2. The system lists the questions with their status and their use (answers, share correct).
  3. The teacher approves, edits or rejects each pending question.
  4. Each edit goes through the same checks as an LLM draft.}
\ucrow{Alternatives}{1a. The teacher generates new questions for a section; they enter the bank as pending.
  4a. An edit fails a check (for example two identical options): it is refused with the reason.}
\ucrow{Post-conditions}{Approved questions can be served in quizzes; every review is traced with what changed.}
\bottomrule
\end{tabularx}
\end{table}

\begin{table}[htbp]
\centering
\caption{Use case \emph{Analyse the course material}}
\label{tab-uc-material}
\small
\begin{tabularx}{\textwidth}{@{}L{2.6cm}Y@{}}
\toprule
\ucrow{Actors}{Teacher (primary); LLM provider (secondary)}
\ucrow{Pre-conditions}{The course has assessed learners.}
\ucrow{Main scenario}{1. The teacher starts the analysis.
  2. The Material Improvement agent loads the evidence and runs the five detectors.
  3. New issues are recorded and the request returns at once.
  4. In the background, the LLM writes a diagnosis and a draft for each issue; quotes are checked.
  5. The teacher reviews each suggestion: accept (possibly after editing), or reject.
  6. After changing the material in Moodle, the teacher marks the suggestion as applied.}
\ucrow{Alternatives}{2a. No detector fires (for example a small class): no suggestion is created.
  4a. A draft fails: the other drafts continue.}
\ucrow{Post-conditions}{Suggestions and reviews are stored and traced; nothing is written to Moodle.}
\bottomrule
\end{tabularx}
\end{table}

\begin{table}[htbp]
\centering
\caption{Use case \emph{Act on an early warning}}
\label{tab-uc-warning}
\small
\begin{tabularx}{\textwidth}{@{}L{2.6cm}Y@{}}
\toprule
\ucrow{Actors}{Teacher (primary); Scheduler (secondary)}
\ucrow{Pre-conditions}{A context sync has scored the students of the course; the teacher teaches the course.}
\ucrow{Main scenario}{1. The teacher opens the class page.
  2. The system lists the students with a medium or high risk, with the top three reasons and the trend over 7 days.
  3. The teacher opens a warning and reads the pre-written nudge, built from the top reason.
  4. The teacher edits the message and sends it.
  5. The nudge appears in the student's feed with one action (a quiz on the weakest key section, or the overdue activity).
  6. Within 7 days, the page shows whether the student came back (an answer, a reading or a tutor question).}
\ucrow{Alternatives}{3a. The teacher marks the warning handled or dismisses it: it is hidden for 14 days, unless its level rises.
  2a. The student is in the baseline arm of a running study: they are never listed.}
\ucrow{Post-conditions}{The follow-up is stored with the score and the reasons it acted on; nothing is sent without the teacher.}
\bottomrule
\end{tabularx}
\end{table}

% =====================================================================
\section{Architecture}
\label{sec-d-arch}
% =====================================================================

\subsection{Global layered architecture}

\Cref{fig-architecture} presents the architecture in five layers. The
\emph{presentation} layer is a React single-page application with four
areas: the course reader, the tutor chat with its live agent steps, the
adaptive quizzes, and the teacher pages. The \emph{gateway} is a NestJS
API that handles authentication, the course reader, the chat sessions,
the quizzes and the teacher routes; it also contains the only Moodle
client. The \emph{agents} layer is a Python service (FastAPI and
LangGraph) that holds every agent, the \texttt{StateStore}, the tracer
and a worker for background jobs. The \emph{data} layer is PostgreSQL with
pgvector, for application data, the learner state, the embeddings and the
traces, and Redis for the job queues. The \emph{external services} are
Moodle, the chat LLM (hosted or local) and the embedding model.

\begin{figure}[htbp]
  \centering
  \reportfig{architecture}
  \caption{Layered architecture of the framework.}
  \label{fig-architecture}
\end{figure}

The gateway sits between the web app and the agents for two reasons.
First, it is the \emph{only} Moodle client: the agents obtain Moodle data
through internal API endpoints, so the Moodle token and the other secrets
stay in one place. The price is one extra hop (agents, then API, then
Moodle). Second, it is the place for authentication, sessions and the
course reader, which are not agent logic. Conversely, all LLM logic lives
in the Python agents: ingestion, retrieval and the agent graphs. LLM and
embedding access sit behind one provider interface, switched by
environment variables. Long jobs go through BullMQ: the API produces jobs
and a Python worker consumes them.

\subsection{Deployment view}

\Cref{fig-deployment} shows the deployment. All services run as
containers of one Docker Compose project on a single host: the web app
(nginx, port 5173), the API (Node, port 3001), the agents (uvicorn, port
8000) and their worker, Moodle (Apache and PHP, port 8080) with its own
PostgreSQL database, the application database (PostgreSQL with pgvector,
port 5432), Redis (port 6379), and a one-shot \texttt{migrate} container
that applies the schema with dbmate. The local models run in Ollama on
the host machine (port 11434), and the hosted models are reached through
the Groq API.

\begin{figure}[htbp]
  \centering
  \reportfig{deployment}
  \caption{Deployment view (Docker Compose).}
  \label{fig-deployment}
\end{figure}

\subsection{Multi-agent architecture}

\Cref{fig-multi-agent} (\cref{chap-contribution}) showed the agents
around the shared state; \cref{tab-agents-roster} gave their roles and
the part played by the LLM. \Cref{tab-agents-io} completes this view
with the inputs and outputs of each agent. Two rules hold for all of
them: every read and write of the learner state goes through the
\texttt{StateStore} with optimistic concurrency, and every decision is
recorded by the tracer in \texttt{agent\_event}.

\begin{table}[htbp]
\centering
\caption{Inputs and outputs of the agents}
\label{tab-agents-io}
\scriptsize
\begin{tabularx}{\textwidth}{@{}L{2.5cm}YY@{}}
\toprule
\textbf{Agent} & \textbf{Inputs} & \textbf{Outputs} \\
\midrule
Orchestrator & message, language & intent (concept, exercise, quiz, course\_info, smalltalk, off\_topic) and route \\
Tutor & question, passages, mastery, context, style posterior & grounded answer with numbered citations, or refusal; chosen explanation style \\
Assessment & section, learner language, question bank & quiz of 5 items (bank first) \\
Diagnosis & each graded answer & updated mastery $P(L)$, misconceptions \\
Recommendation & mastery, context, objectives, learner goal, teacher overrides & ranked sections with structured reasons; learning path \\
Context & Moodle sources, in-app views & course and learner context, interventions \\
Learning Analytics & all learner states of the course, recent activity & per-section and per-student view, risk flags; risk score, level and reasons, daily snapshot \\
Material Improvement & answers, unanswered turns, reading, mastery, views & suggestions with evidence and drafts \\
Ingestion (worker) & Moodle contents and files & chunks and embeddings \\
\bottomrule
\end{tabularx}
\end{table}

\subsection{Main design decisions}

\Cref{tab-adr} summarizes the twelve architecture decision records, with
the alternatives considered and the reason for each choice.

\begin{table}[htbp]
\centering
\caption{Summary of the architecture decision records}
\label{tab-adr}
\scriptsize
\begin{tabularx}{\textwidth}{@{}L{0.9cm}YL{3.2cm}L{3.3cm}@{}}
\toprule
\textbf{ADR} & \textbf{Decision} & \textbf{Alternatives} & \textbf{Reason} \\
\midrule
0001 & Layered architecture; Moodle as backend only; NestJS is the only Moodle client; Python agents hold all LLM logic; BullMQ for long jobs & agents calling Moodle directly; a Moodle plugin & one place for the Moodle token and secrets; core untouched \\
0002 & One shared, versioned learner state in PostgreSQL, written through the StateStore with optimistic concurrency; a snapshot per turn; \texttt{agent\_event} trace & private agent memories & consistency between agents, auditability \\
0003 & Adaptive loop with BKT per section; quizzes validated in code; recommendation rules & IRT, DKT, reinforcement learning & interpretable, works with few answers, defensible \\
0004 & Explicit context model, Context agent and rule-based interventions, polling Moodle web services & a Moodle event plugin; LLM-decided triggers & no core change; explainable triggers \\
0005 & Teacher in the loop: question bank with review, objectives, class dashboard & auto-publish only & a measured faithfulness of about 0.5 requires human review \\
0006 & In-app course reader; signed file links (HMAC, 1\,h); in-app views recorded, never written to Moodle & links to Moodle; writing completion back & one interface, no double login; Moodle stays read-only \\
0007 & Material Improvement: code detects (statistics), the LLM drafts, the teacher validates; on demand & LLM-only critique; scheduled runs & explainable, measurable, quota-friendly \\
0008 & Agent steps in the chat from the graph's own spans, never model reasoning; course-centred interface & showing the LLM's chain of thought & transparency without leaking solutions \\
0009 & Learning-gain evaluation: simulated ablation (full, without context, baseline) and a two-arm pilot kit; prerequisites of a close-deadline section first & a pilot only; deadline before prerequisite & an evaluation now, and a ranking rule it showed to be better \\
0010 & Learner goal, goal-aware path, structured reasons, explanation-style bandit, teacher pin and skip & LLM-written explanations; full RL agent; learning-styles questionnaire & real reasons; learns from a few ratings; weak evidence for learning styles \\
0011 & LTI 1.3 launch; progress grades through AGS; consent, pseudonyms and analysis for the pilot & dev login; grade web services; reading Moodle logs & no second login, no new capability, core untouched \\
0012 & Logistic risk model fitted on simulated classes, with reasons; the teacher decides each follow-up & hand-weighted score; gradient boosting with SHAP; automatic nudges & honest reasons, an evaluated warning, teacher in control \\
\bottomrule
\end{tabularx}
\end{table}

% =====================================================================
\section{Dynamic view}
\label{sec-d-dynamic}
% =====================================================================

This section describes the behaviour of the system in its four main
scenarios with sequence diagrams, then two scenarios added in the later
sprints (the launch from Moodle and the early warning), the ingestion of
a course with an activity diagram, and the life cycle of three objects
with state machines.

\subsection{A tutor turn}

\Cref{fig-seq-tutor} shows a tutor turn. The student's question reaches
the API, which forwards it to the agents and proxies their stream back to
the web app. The agents first load the learner state and its version
from the \texttt{StateStore}. The Orchestrator classifies the question,
and its step is streamed to the chat at once. The Tutor retrieves the
closest passages and streams a second step. If the best similarity
reaches the threshold, the Tutor adapts the level to the learner's
mastery and context, asks the LLM to write the answer, streams the tokens
as they arrive, and checks the citations; otherwise it streams a
\texttt{not\_found} step with the closest sections. In both cases,
\texttt{commit\_turn} writes the messages, the steps and the new context
snapshot in one transaction, and the final events (citations, end of the
turn) close the stream.

\begin{figure}[htbp]
  \centering
  \reportfig{seq-tutor}
  \caption{Sequence of a tutor turn with streamed agent steps.}
  \label{fig-seq-tutor}
\end{figure}

\subsection{An adaptive quiz}

\Cref{fig-seq-quiz} shows an adaptive quiz. The quiz is created at once
in the state \emph{generating} and returned to the web app, which polls
it; the work continues in the background. The Assessment agent first takes
approved bank questions that the learner has never answered. If some are
missing, the review mode decides: in \texttt{auto\_publish} mode the LLM
writes the missing items, which are validated in code and inserted as
approved; in \texttt{review} mode they are inserted as pending, and the
quiz is served only if at least $n-2$ questions are approved, otherwise it
ends with \texttt{awaiting\_review}. For each answer, the grading is done
in code, and the BKT update is committed atomically with the response;
the reply carries the mastery before and after. Finally, the
Recommendation agent ranks the sections and gives the reason of each.

\begin{figure}[htbp]
  \centering
  \reportfig{seq-quiz}
  \caption{Sequence of an adaptive quiz.}
  \label{fig-seq-quiz}
\end{figure}

\subsection{A context sync and proactive interventions}

\Cref{fig-seq-context} shows a context sync. It starts in one of two
ways: the repeatable job of the scheduler, every 10 minutes, or a learner
who opens a course, which enqueues a sync at most once a minute. The
worker consumes the job and asks the API for the context sources of the
course; the API reads them from Moodle and adds the in-app views. The
Context agent builds and commits the course context, then, for each
enrolled learner, builds the learner context, commits it with a version
check, evaluates the four triggers and inserts the new interventions,
using their de-duplication key. Finally, the Learning Analytics agent
scores the risk of every student of the course and upserts the
snapshot of the day (\cref{sec-c12}). Later, the student's feed shows
each suggestion with its reason.

\begin{figure}[htbp]
  \centering
  \reportfig{seq-context}
  \caption{Sequence of a context sync and proactive interventions.}
  \label{fig-seq-context}
\end{figure}

\subsection{A material analysis}

\Cref{fig-seq-material} shows a material analysis. The teacher's request
reaches the Material service in the agents. The service loads the
evidence, embeds the unanswered questions to group them, runs the five
detectors and records the issues with their stable keys; the request then
returns with the status \texttt{202} (accepted). The drafts are written
in the background, one after the other: for each new issue the LLM writes
a diagnosis and a draft, the quote is checked, and the suggestion is saved
as pending. The web app polls until the drafts are ready. The teacher then
accepts, edits or rejects each suggestion, and marks it as applied after
changing the material in Moodle; each change is traced.

\begin{figure}[htbp]
  \centering
  \reportfig{seq-material}
  \caption{Sequence of a material analysis.}
  \label{fig-seq-material}
\end{figure}

\subsection{A launch from Moodle}

A student opens the ``Adaptive tutor'' activity of the course in Moodle.
Moodle sends an OpenID Connect login request to \texttt{/lti/login}; the
API checks the issuer and the client, stores a single-use \texttt{state}
and \texttt{nonce} in Redis and redirects the browser to Moodle's
authorization endpoint. Moodle then posts a signed \texttt{id\_token} to
\texttt{/lti/launch}. The API verifies the token against Moodle's key set
and checks its claims, records the deployment and the resource link with
its line item, creates or updates the user, the course and the
enrolment, and redirects the browser to the web app with the app's own
session token in the URL fragment. The web app reads the token, removes
it from the address bar and opens the course. The API then posts the
student's current progress to the line item, as it does after every quiz
answer.

\subsection{An early warning and a teacher nudge}

After each context sync, the Learning Analytics agent computes the
eleven risk features of each student from the shared state and the
recent activity, applies the logistic model, keeps the three largest
positive contributions as reasons, and upserts the snapshot of the day.
When the teacher opens the class page, the API asks the agents for the
class view, which scores the students live and reads the snapshot of
7 days earlier for the trend. A nudge sent by the teacher is stored as a
follow-up with the score and the reasons, and as a \texttt{teacher\_nudge}
intervention in the student's feed. The next views of the class page
check the student's activity since the nudge and show whether they came
back.

\subsection{Course ingestion}

\Cref{fig-ingestion} shows the ingestion of a course as an activity
diagram with two swimlanes. The teacher starts the indexing, and the API
enqueues a job. The worker fetches the course contents through an
internal API endpoint and processes each activity: it downloads the
content and skips it if its hash has not changed; otherwise it extracts
the text (HTML or PDF), detects its language, cuts it into chunks,
computes the embeddings with bge-m3, and upserts the document and its
chunks, with their vectors and their full-text index. When every activity
is processed, the documents of activities removed from the course are
deleted, and the job ends with its statistics. Because unchanged content
is skipped and documents are upserted, a new ingestion of the same course
is idempotent.

\begin{figure}[htbp]
  \centering
  \reportfig{ingestion-activity}
  \caption{Course ingestion pipeline.}
  \label{fig-ingestion}
\end{figure}

\subsection{State machines}

\Cref{fig-states} gives the life cycles of three objects. A \emph{quiz}
is generating, then ready, then completed when all its items are
answered; it fails with a code when the LLM fails, times out or when the
quiz waits for a teacher review. A \emph{bank question} generated in
\texttt{review} mode is pending until the teacher approves or rejects it;
in \texttt{auto\_publish} mode it is approved at once, and an edit sends
it back to review. A \emph{material suggestion} is pending, then accepted
or rejected, and an accepted suggestion becomes applied; each of these
transitions can be undone.

\begin{figure}[htbp]
  \centering
  \reportfig{states}
  \caption{State machines of quizzes, questions and material suggestions.}
  \label{fig-states}
\end{figure}

% =====================================================================
\section{Static view}
\label{sec-d-static}
% =====================================================================

\Cref{fig-class-diagram} presents the domain classes, grouped into six
packages.

\begin{figure}[htbp]
  \centering
  \reportfig{class-diagram}
  \caption{Domain class diagram (core).}
  \label{fig-class-diagram}
\end{figure}

\begin{description}[style=nextline, leftmargin=0.6cm]
  \item[Identity] \texttt{User} and \texttt{Course} mirror the Moodle user
    and course, with their Moodle identifiers; \texttt{Enrolment} holds
    the role of a user in a course and the last access recorded by the
    app.
  \item[Learner state] \texttt{LearnerCourseState} is the shared learner
    state of one learner in one course, with its \texttt{version}. It
    maps each topic to a \texttt{TopicMastery} (the probability
    \texttt{p\_known} and the number of attempts), and contains the
    misconceptions, the interaction statistics and the
    \texttt{LearnerContext}.
  \item[Context] \texttt{CourseContext} holds the deadlines and the
    prerequisite graph of a course, with its own version.
    \texttt{Intervention} is a proactive suggestion for a learner in a
    course, with its trigger, reason, action, de-duplication key and
    status. \texttt{ActivityView} records the views and reading time of an
    activity in the app, which feed the learner context.
  \item[Assessment] A \texttt{Quiz} on one topic contains
    \texttt{QuizItem}s; each item may be a copy of a bank
    \texttt{Question}. A \texttt{QuizResponse} records the chosen option,
    whether it was correct, and the mastery before and after.
  \item[Teaching] The \texttt{CoursePolicy} of a course holds its review
    mode and its objectives. A \texttt{MaterialSuggestion} holds a
    proposed change with its evidence, diagnosis, draft and status.
  \item[Chat and tracing] A \texttt{LearningSession} contains
    \texttt{Message}s; an assistant message keeps its citations, intent,
    trace identifier and its \texttt{AgentStep}s.
\end{description}
Four groups of classes were added in the later sprints; \cref{fig-class-diagram-ext}
shows them around the core classes they refer to.

\begin{figure}[htbp]
  \centering
  \reportfig{class-diagram-ext}
  \caption{Classes added in sprints 8 to 12: personalization, study, LTI and early warning.}
  \label{fig-class-diagram-ext}
\end{figure}

\begin{description}[style=nextline, leftmargin=0.6cm]
  \item[Personalization] The \texttt{Personalization} part of
    \texttt{LearnerCourseState} holds the goal, the level, the preferred
    style and the Beta posterior of each explanation style.
    \texttt{MessageFeedback} records a rating of an answer, and
    \texttt{RecommendationOverride} a pin or a skip by the teacher.
  \item[Study] \texttt{StudyCondition} gives the arm and the pseudonym of
    a student in a study; \texttt{StudyConsent} their answer and any
    withdrawal. A \texttt{KnowledgeTest} (pre or post) holds copied items,
    a \texttt{KnowledgeTestAttempt} the answers and the score, and a
    \texttt{SusResponse} the questionnaire and its score.
  \item[LTI] \texttt{LtiDeployment} records a platform deployment, and
    \texttt{LtiResourceLink} the activity of a course with its line item.
  \item[Early warning] A \texttt{RiskSnapshot} holds the features, the
    score, the level and the reasons of a student for one day, with the
    model version; a \texttt{RiskFollowup} records what the teacher did
    and the score and reasons it acted on.
\end{description}
The \guillemotleft service\guillemotright{} \texttt{StateStore} is the only
way for the agents to load and commit a learner state or a course
context; its \texttt{commit\_turn} operation applies the optimistic
concurrency described in \cref{sec-c1}.

% =====================================================================
\section{Database design}
\label{sec-d-db}
% =====================================================================

\Cref{fig-database} shows the relational schema of the core tables, \cref{fig-database-ext}
the tables added in sprints 8 to 12, and \cref{tab-db-tables}
describes the role of each table. The schema is owned by SQL migrations
applied by one tool, dbmate. The API reads it through types generated for
Kysely, and the agents through Pydantic models that mirror the tables.

\begin{figure}[htbp]
  \centering
  \reportfig{database}
  \caption{Relational schema of the database (core tables).}
  \label{fig-database}
\end{figure}

\begin{figure}[htbp]
  \centering
  \reportfig{database-ext}
  \caption{Tables added in sprints 8 to 12.}
  \label{fig-database-ext}
\end{figure}

\begin{table}[htbp]
\centering
\caption{Tables of the database}
\label{tab-db-tables}
\scriptsize
\begin{tabularx}{\textwidth}{@{}L{2.0cm}L{3.4cm}YL{3.6cm}@{}}
\toprule
\textbf{Domain} & \textbf{Table} & \textbf{Role} & \textbf{Key columns} \\
\midrule
Identity & \texttt{app\_user} & mirror of a Moodle user & moodle\_user\_id, role, locale \\
 & \texttt{course} & mirror of a Moodle course & moodle\_course\_id, shortname \\
 & \texttt{enrolment} & role of a user in a course & role, last\_access\_at \\
Content / RAG & \texttt{ingestion\_job} & one indexing job & status, stats \\
 & \texttt{course\_document} & one ingested source (activity or file) & moodle\_cmid, section\_num, lang, content\_hash \\
 & \texttt{course\_chunk} & a chunk with its vector and full-text index & embedding vector(1024), tsv \\
Learner state & \texttt{learner\_profile} & cross-course learner facts & preferred language, version \\
 & \texttt{learner\_course\_state} & the shared learner state & mastery, misconceptions, interaction\_stats, context (jsonb), version \\
 & \texttt{activity\_view} & views and reading time in the app & cmid, views, seconds \\
Chat \& tracing & \texttt{learning\_session} & a conversation & user, course \\
 & \texttt{message} & a message, with citations and steps & content, citations, intent, steps, trace\_id \\
 & \texttt{context\_snapshot} & the state seen and produced at each turn & state (jsonb), version \\
 & \texttt{agent\_event} & audit log of decisions and steps & trace\_id, agent, kind, payload, latency\_ms \\
Assessment & \texttt{quiz} & a quiz on one topic & topic\_key, lang, status \\
 & \texttt{quiz\_item} & an item, possibly copied from the bank & stem, options, answer\_index, question\_id \\
 & \texttt{quiz\_response} & an answer and its BKT effect & chosen\_index, correct, p\_known\_before, p\_known\_after \\
Context & \texttt{course\_context} & the course context & context (jsonb), version \\
 & \texttt{intervention} & a proactive suggestion & trigger, dedupe\_key, reason, action, status \\
Teacher & \texttt{question} & a reviewed bank question & status, options, answer\_index, reviewer \\
 & \texttt{course\_policy} & review mode and objectives & review\_mode, objectives (jsonb), version \\
 & \texttt{material\_suggestion} & a proposed change to the material & kind, evidence, diagnosis, draft, status \\
Personalization & \texttt{message\_feedback} & a rating of a tutor answer & message, rating, reason \\
 & \texttt{recommendation\_override} & a pin or skip by the teacher & topic\_key, action, note, user (or class) \\
Study & \texttt{study\_condition} & arm of a student in a study & arm, pseudonym \\
 & \texttt{study\_consent} & consent and withdrawal & consented, withdrawn\_at \\
 & \texttt{knowledge\_test} & a pre-test or post-test & phase, items (jsonb) \\
 & \texttt{knowledge\_test\_attempt} & one student's test & answers, correct, total \\
 & \texttt{sus\_response} & the SUS questionnaire & answers, score \\
LTI & \texttt{lti\_deployment} & a platform deployment & issuer, client\_id, deployment\_id \\
 & \texttt{lti\_resource\_link} & an LTI activity of a course & resource\_link\_id, lineitem\_url \\
Early warning & \texttt{risk\_snapshot} & daily risk of a student & day, score, level, features, reasons, model\_version \\
 & \texttt{risk\_followup} & a teacher's follow-up & kind, message, score, reasons \\
\bottomrule
\end{tabularx}
\end{table}

Three choices shape this design:
\begin{itemize}
  \item \emph{One database for everything.} PostgreSQL with the pgvector
        extension holds the relational data, the learner state as
        \texttt{jsonb} and the vectors of the course chunks. Retrieval can
        therefore combine vector search and full-text search in one query
        engine, and there is no second store to keep consistent.
  \item \emph{\texttt{jsonb} plus a version column for shared state.} The
        parts of the learner state, the course context and the course
        policy are \texttt{jsonb} documents, so they can evolve without a
        migration, and each row carries a \texttt{version} that implements
        the optimistic concurrency of the \texttt{StateStore}.
  \item \emph{Migrations owned by one tool.} The schema is defined by SQL
        migrations applied by dbmate in a one-shot container, and both the
        API and the agents derive their types from it.
\end{itemize}

% =====================================================================
\section{Technology stack}
\label{sec-d-stack}
% =====================================================================

\Cref{fig-stack} summarizes the technologies by layer, and
\cref{tab-stack} gives their pinned versions and the reason for each
choice.

\begin{figure}[htbp]
  \centering
  \reportfig{stack}
  \caption{Technology stack.}
  \label{fig-stack}
\end{figure}

\begin{table}[htbp]
\centering
\caption{Technology stack, versions and reasons}
\label{tab-stack}
\scriptsize
\begin{tabularx}{\textwidth}{@{}L{1.8cm}L{3.6cm}L{1.8cm}Y@{}}
\toprule
\textbf{Layer} & \textbf{Technology} & \textbf{Version} & \textbf{Why} \\
\midrule
Presentation & React, TypeScript, Vite & 19.3, 6.0, 8 & own interface, independent of Moodle's theme \\
 & Tailwind CSS, shadcn/ui & 4.3, -- & customisable, modern components \\
 & TanStack Query & 5 & server state and polling (quizzes, drafts) \\
 & react-i18next & -- & French and English interface \\
Gateway & NestJS (Node $\geq$ 22) & 12.3 & structured backend for auth, sessions and the Moodle client \\
 & Kysely, zod & -- & typed SQL from the schema; validation of requests \\
 & BullMQ (producer), JWT & 6.5, -- & background jobs; authentication \\
 & jose & -- & LTI 1.3: verification of Moodle's tokens, signed assertions for AGS \\
Agents & Python (uv), FastAPI & 3.12, 0.162 & Python ecosystem for LLMs, retrieval and evaluation \\
 & LangGraph, LangChain & 1.2, 1.4 & explicit agent graphs with conditional edges and custom streams \\
 & pydantic, psycopg & --, 3 & models of the shared state; PostgreSQL access \\
 & bullmq (consumer) & 3.4 & consumes the jobs produced by the API \\
 & Ragas & 0.4 & evaluation of faithfulness \\
 & SciPy (development only) & 1.18 & statistics of the pilot analysis \\
LLMs & Groq: \texttt{qwen/qwen3.8-27b} (tutor, router), \texttt{openai/gpt-oss-20b} (quiz and material writer) & -- & hosted models for quality and speed \\
 & Ollama: \texttt{qwen3:4b-instruct}, \texttt{bge-m3} (1024 dimensions) & -- & local option for privacy; multilingual embeddings \\
Data & PostgreSQL + pgvector & 16 & one database for relational data, JSONB state and vectors \\
 & Redis & 7 & BullMQ queues \\
LMS & Moodle LTS, image \texttt{moodlehq/} \texttt{moodle-php-apache} 8.3 & 4.5.15 & long-term support; core never modified \\
Infrastructure & Docker Compose, dbmate, \texttt{just} & -- & reproducible stack, one migration tool, task runner \\
Quality & Vitest, msw, pytest, ESLint, Prettier, ruff, mypy & -- & tests (LLMs and Moodle replaced by fakes) and static checks \\
\bottomrule
\end{tabularx}
\end{table}

The design above specifies what the platform must do and how it is
built: the requirements and use cases, the layered architecture with its
gateway and its agents around one shared state, the behaviour of the main
scenarios, the domain model, the database and the technology stack, each
choice traced back to an architecture decision record. The following
sections present the development of the platform.
